\subsection{Fréchet 空间的严格态射}

定理 1. --- 设 E、F 是两个 Fréchet 空间，u 是 E 到 F 中的一个连续线性映射。下列条件等价：

\begin{enumerate}
\def\labelenumi{\alph{enumi})}
\item
  \(u\) 是一个严格态射。
\item
  \(u\) 对于弱化拓扑是一个严格态射。
\item
  u 的像在 F 中是闭的。
\item
  \({}^t u\) 对于弱拓扑是 \({\mathrm{F}}^{\prime }\) 到 \({\mathrm{E}}^{\prime }\) 中的一个严格态射。
\item
  \({}^t u\) 的像在 \({\mathrm{E}}^{\prime }\) 中对于弱拓扑 \(\sigma \left( {{\mathrm{E}}^{\prime },\mathrm{E}}\right)\) 是闭的。
\item
  \({}^t u\) 的像对于强拓扑 β(E', E) 在 E' 中是闭的。
\end{enumerate}

\(g)\) \({}^t u\) 是 \(\mathrm{F}_c^\prime\) 到 \(\mathrm{E}_c^\prime\) 中的一个严格态射（两个对偶空间均赋予紧收敛拓扑）。

\(a\))、\(b\)) 与 \(e\)) 的等价性由 IV 第 28 页的命题 3 及其前面的注记推出；\(a\)) 与 \(c\)) 的等价性恰好就是 I 第 19 页的推论 3。IV 第 28 页的注记还表明，\(d\)) 等价于下述事实：\(u\) 的像对于 \(\sigma(F, F')\)（即 \(F\) 的弱化拓扑）是闭的；\(c\)) 与 \(d\)) 的等价性随即由 IV 第 4 页的命题 2 推出。

我们现在来证明 \(e)\) 与 \(f)\) 的等价性。只需证明 \(f)\) 蕴含 \(e)\)。设 \({}^t u\) 的像对于 \(\beta(E', E)\) 而言在 \(E'\) 中是闭的。根据 Banach-Dieudonné 定理 (IV, 第 25 页, 推论 2)，只需证明：对于每个凸均衡邻域 \(U\)（即 \(0\) 在 \(E\) 中的一个邻域），交 \(B = {}^t u(F') \cap U\) 对于 \(\sigma(E', E)\) 是紧的。\(E_b'\)（Fréchet 空间 \(E\) 的强对偶）是完备的 (IV, 第 22 页, 命题 2)，因此闭子集 \(B\)（\(E_b'\) 中的）是完备的，从而赋范空间 \(E_B'\) 是完备的 (III, 第 8 页, 推论)。设 \((V_n)\) 是一个递减序列，构成 \(0\) 在 \(F\) 中的基本邻域系。那么 \(F'\) 是诸集 \(C = V_n^\circ\) 之并，而这些集对于 \(\sigma(F', F)\) 都是紧的，于是 \(E_B' = \bigcup B_n\)，其中 \(B_n = E_B' \cap {}^t u(C_n)\)。由于 \(E_B'\) 是一个 Baire 空间，且诸集 \(B_{n}\) 都是凸的、均衡的且闭的，故存在实数 r \textgreater{} 0 与整数 n，使得 \(B \subset r \cdot B_{n}\)。这时我们有 \(\mathbf{B} = \mathbf{U}^{\circ} \cap {}^{t} u(r \cdot C_{n})\)；由于诸集 \(U^{\circ}\) 与 \(r \cdot C_{n}\) 都是紧集，而 \({}^{t} u\) 对于弱拓扑是连续的，所以 B 对于 \(\sigma(\mathrm{E}', \mathrm{E})\) 是紧的。这就完成了 e) 与 f) 等价性的证明。

最后，\(g\)) 与前面诸条件的等价性由 GT, X, § 2, 第 10 目的命题 18 以及下面的引理推出。

引理 1. --- 设 E、F 是两个 Hausdorff 局部凸且拟完备的空间，u 是 E 到 F 中的一个连续线性映射。要使 \({}^t u\) 是 \(F_{c}'\) 到 \(E_{c}'\) 中的一个严格态射，必须且只需 u 的像 \(u(E)\) 是闭的，并且 \(u(E)\) 的每个紧子集都是 E 的某个紧子集在 u 下的像。

由 Mackey 定理 (IV, 第 2 页, 定理 1) 以及在 \(\mathrm{E}'\)（相应地，F）上紧收敛拓扑与紧凸收敛拓扑相重合这一事实 (IV, 第 4 页)，我们可以把 E（相应地，F）与对偶 \(\mathrm{E}_c'\)（相应地，\(\mathrm{F}_c'\)）等同起来。于是 \(u\) 是 \({}^t u\) 的转置，而 E（相应地，F）的等度连续子集就是相对紧集。引理 1 于是由 (IV, 第 27 页) 的命题 2 推出，因为 \(u(\mathrm{E})\) 在 F 中是闭的当且仅当它对于弱化拓扑 \(\sigma(\mathrm{F}, \mathrm{F}')\) 是闭的 (IV, 第 4 页, 命题 2)。

推论 1. --- 在定理 1 的假设下，下列条件等价：

\begin{enumerate}
\def\labelenumi{(\roman{enumi})}
\item
  \(u\) 是一个严格单射态射；
\item
  \({}^{t}u\) 对于弱拓扑是一个严格满射态射。
\item
\({}^{t}u\) 是满射。
\end{enumerate}

(i) \(\Rightarrow\) (ii) 可由定理 1 中条件 \(a\))、\(d\)) 与 \(e\)) 的等价性以及 IV 第 6 页的命题 5 立即得出。显然 (ii) 蕴含 (iii)。最后证明 (iii) 蕴含 (i)：若 \({}^t u\) 是满射，则 \(u\) 由定理 1 中 \(a\)) 与 \(e\)) 的等价性可知是一个严格态射；而 \(u\) 是单射这一点由 IV 第 6 页的命题 5 得出。

推论 2. --- 在定理 1 的假设下，下列条件等价：

\begin{enumerate}
\def\labelenumi{(\roman{enumi})}
\item
  \(u\) 是满射；
\item
  \(u\) 是一个严格满射态射；
\item
  \({}^t u\) 对于弱拓扑是一个严格单射态射。
\end{enumerate}

(i) 与 (ii) 的等价性由 Banach 定理 (I, 第 17 页, 定理 1) 得出。

鉴于定理 1 中 \(a\) 与 \(c\) 的等价性，条件 (ii) 表明 \(u\) 是一个严格态射，并且它的像在 F 中对于 \(\sigma(F, F')\) 是稠密的。(ii) 与 (iii) 的等价性随后由定理 1 中 \(a\) 与 \(d\) 的等价性以及 IV 第 6 页的命题 5 得出。

若 \(u: \mathrm{E} \to \mathrm{F}\) 是 Fréchet 空间之间的一个严格态射，则转置 \({}^t u\) 不一定是 \(\mathrm{F}_b'\) 到 \(\mathrm{E}_b'\) 中的严格态射 (IV, 第 62 页, 习题 3)。尽管如此，我们有下面的部分结果：

推论 3. --- 在定理 1 的假设下，下面的性质蕴含性质 a) 至 g)：

\begin{enumerate}
\def\labelenumi{\alph{enumi})}
\setcounter{enumi}{7}
\tightlist
\item
  \({}^{t}u\) 是 \(F_{b}^{\prime}\) 到 \(E_{b}^{\prime}\) 中的一个严格态射。
\end{enumerate}

当 E 与 F 都是 Banach 空间或都是 Montel 空间时，性质 \(h)\) 等价于定理 1 的性质 \(a)\) 至 \(g)\)。

设 \({}^t u\) 是 \(\mathrm{F}_b'\) 到 \(\mathrm{E}_b'\) 中的一个严格态射。我们将证明 \({}^t u\) 的像 H 在 \(\mathrm{E}_b'\) 中是闭的，由此即得推论 3 的第一个断言。

设 G 是 u 的像在 F 中的闭包；空间 G 赋予由 F 诱导的拓扑后是一个 Fréchet 空间。映射 \(u: E \to F\) 分解为 \(u = j \circ v\)，其中 j 是 G 到 F 中的典范单射，而 \(v \in \mathcal{L}(E; G)\)。于是有 \({}^t u = {}^t v \circ {}^t j\)，其中 \({}^t j\) 是满射（由 Hahn-Banach 定理, II, 第 24 页, 命题 2）；又，\({}^t v\) 是单射，因为 \(v(E)\) 在 G 中稠密 (IV, 第 6 页, 命题 5)。根据假设，映射 \({}^t u\) 从 \(F_{b}'\) 到 H 上是开的；由于 \({}^t j\) 是满射且连续，映射 \({}^t v\) 诱导出 \(G_{b}'\) 到 H 上的一个同胚。而 Fréchet 空间 G 的对偶 \(G_{b}'\) 是完备的 (IV, 第 22 页, 命题 2)；因此 H 是完备的，从而在 \(E_{b}'\) 中是闭的。

若 E 与 F 是 Montel 空间，则 \(E'\)（相应地，\(F'\)）上的强拓扑与紧收敛拓扑重合，而 h) 只是 g) 的一个重新表述。

若 E 与 F 是 Banach 空间，则 \(E_{b}^{\prime}\) 与 \(F_{b}^{\prime}\) 也是 Banach 空间，而条件 h) 等价于 f)，这是把 a) 与 c) 的等价性应用于 \(u: F_{b}^{\prime} \to E_{b}^{\prime}\) 所得的结果。

推论 4. --- 设 E 与 F 是 Banach 空间。要使 \({}^t u\) 是满射，必须且只需存在实数 \(r > 0\)，使得 \(\| x\| \leqslant r.\| u(x)\|\) 对所有 \(x\in E\) 成立。

这只是推论 1 中条件 (i) 与 (iii) 等价性的一个重新表述。

推论 5. --- 设 E、F 是两个 Fréchet 空间，u 是一个连续线性映射，

它从 E 映到 F 中。下列条件等价：

\begin{enumerate}
\def\labelenumi{\alph{enumi})}
\item
  \(u\) 是 E 到 F 上的一个同构。
\item
  u 对于弱化拓扑是 E 到 F 上的一个同构。
\item
  \({}^t u\) 对于弱拓扑是 \({\mathrm{F}}^{\prime }\) 到 \({\mathrm{E}}^{\prime }\) 上的一个同构。
\item
  \({}^{t}u\) 对于强拓扑是 \(F'\) 到 \(E'\) 上的一个同构。
\item
  \({}^t u\) 是 \({\mathrm{F}}_{c}^{\prime }\) 到 \({\mathrm{E}}_{c}^{\prime }\) 上的一个同构。
\end{enumerate}

由于同构无非就是双射的严格态射，\(a)\) 与 \(b)\) 的等价性由定理 1 中条件 \(a)\) 与 \(b)\) 的等价性得出。

显然 a) 蕴含条件 c)、d) 与 e) 中的每一个。

反之，设条件 \(c\))、\(d\) 或 \(e\)) 之一成立。由定理 1 及其推论 3 可知 \(u\) 是 E 到 F 中的一个严格态射，且 \({}^t u\) 显然是双射。设 N（相应地，I）为 \(u\) 的核（相应地，像）。由于 \({}^t u\) 是双射，我们有 \(\operatorname{Im} {}^t u = \operatorname{E}'\) 与 \(\operatorname{Ker} {}^t u = \{0\}\)，于是由 IV 第 27 页的命题 2 得 \(\mathrm{N}^\circ = \mathrm{E}'\) 与 \(\mathrm{I}^\circ = \{0\}\)。但 N（相应地，I）是 E（相应地，F）的一个闭向量子空间，而双极定理 (II, 第 44 页) 蕴含 \(\mathrm{N} = \{0\}\) 与 \(\mathrm{I} = \mathrm{F}\)，故 \(u\) 是双射。于是我们证明了 \(u\) 是一个同构。

